% META: {"id":"1SPE-EXPO-EX-048","chapitre":"1SPE-EXPONENTIELLE","type_objet":"exercice","parcours":3,"competences":["raisonner","calculer"],"duree_min":15,"mode_creation":"ex_nihilo","sources_inspiration":[],"fichier_tex":"chapitres/1SPE-EXPONENTIELLE/exercices/1SPE-EXPO-EX-048.tex","corrige_tex":"chapitres/1SPE-EXPONENTIELLE/corriges/1SPE-EXPO-CO-048.tex","status":"approved","extension_codes":["X1"],"programme_alignment":"OPTIONAL_EXTENSION","extension_label":"Approfondissement — Vers la Terminale","origin":"generated","status_history":["generated","audited","accepted","approved"],"release_acceptance":"RELEASE_OWNER_BATCH_ACCEPTANCE","acceptance_closure_digest":"sha256:9b3ccf9a81c5520fbb7e03b7b2d3e7bf2057a02834b6d908bf3be5f49f3b553f","acceptance_receipt_digest":"sha256:4fad86f0282e0fa4e0419cca1ad6379ae7af5a280c04a4a90880f90a1c044242"}
\begin{center}\textbf{Approfondissement — Vers la Terminale}\end{center}
% BEGIN-VERIFY
% from sympy import *
% x = symbols('x', real=True)
% f = exp(x) / (exp(x) + 1)
% fp = simplify(diff(f, x))
% # f'(x) = e^x / (e^x + 1)^2
% assert simplify(fp - exp(x)/(exp(x) + 1)**2) == 0
% # f'(x) > 0 pour tout x, donc f strictement croissante
% # limites
% assert limit(f, x, oo) == 1
% assert limit(f, x, -oo) == 0
% # f(0) = 1/2
% assert f.subs(x, 0) == Rational(1, 2)
% END-VERIFY
\begin{exercice}{1SPE-EXPO-EX-048}{3}{15}
On considère la fonction $f$ définie sur $\mathbb{R}$ par $f(x) = \dfrac{\mathrm{e}^x}{\mathrm{e}^x + 1}$.
\begin{enumerate}
\item Montrer que $f(x) = \dfrac{1}{1 + \mathrm{e}^{-x}}$ pour tout réel $x$.
\item Calculer $f'(x)$ et montrer que $f'(x) = \dfrac{\mathrm{e}^x}{(\mathrm{e}^x + 1)^2}$.
\item En déduire que $f$ est strictement croissante sur $\mathbb{R}$.
\item Déterminer les limites de $f$ en $+\infty$ et en $-\infty$. En déduire les asymptotes.
\item Calculer $f(0)$ et dresser le tableau de variations complet de $f$.
\end{enumerate}
\end{exercice}
